\section{Three integrals for one number}
\label{sec:vsthree}

We have one number, $w_{\rm vs}=-2\,\E[X]$, defined by the terminal
law.  But the three model families of Part~I hold that law in three
different coordinate systems --- three \emph{charts}, in the word
Chapter~3 introduced for a choice of coordinates: Chapter~3's families
state the \emph{smile} $w(k)$, Chapter~2's model states the \emph{law}
itself through its quantile, and Chapter~4's model states the
\emph{field} $v(\tau,y)$ that generates every law at once.  This
section derives one exact integral for $w_{\rm vs}$ in each chart.
The three integrals look nothing alike, and watching them return the
same number --- to a tenth of a basis point --- is the section's
payoff.

\subsection{Against option prices}

The first route rewrites the expectation of \emph{any} smooth payoff as
a portfolio of the options themselves.

\begin{proposition}[Spanning]
\label{prop:vsspanning}
Let $\varphi$ be twice continuously differentiable on $(0,\infty)$ and
let $Y>0$ have $\E[Y]=1$.  Write $c$ and $P$ for the normalized call
and put at strike $K$, i.e.\ $c=\E[(Y-K)^+]$ and $P=\E[(K-Y)^+]$.
Then, whenever the integrals converge,
\begin{equation}
\E[\varphi(Y)]
=\varphi(1)
+\int_0^{1}\varphi''(K)\,P(K)\,\dd K
+\int_1^{\infty}\varphi''(K)\,c(K)\,\dd K .
\label{eq:vsspanning}
\end{equation}
\end{proposition}

\begin{proof}
The claim reduces to a pointwise identity: for every $y>0$,
\begin{equation}
\varphi(y)=\varphi(1)+\varphi'(1)(y-1)
+\int_0^{1}\varphi''(K)(K-y)^+\dd K
+\int_1^{\infty}\varphi''(K)(y-K)^+\dd K .
\label{eq:vspointwise}
\end{equation}
Check it for $y\ge1$; the case $y<1$ is the mirror image.  For
$y\ge1$ the put-type kernel $(K-y)^+$ vanishes on $K\le1$, and the
call-type kernel is nonzero only for $K\in[1,y]$, so the right-hand
side's integrals reduce to $\int_1^{y}\varphi''(K)(y-K)\,\dd K$.
Integrate by parts once, differentiating $(y-K)$:
\[
\int_1^{y}\varphi''(K)(y-K)\,\dd K
=\Big[\varphi'(K)(y-K)\Big]_1^{y}
+\int_1^{y}\varphi'(K)\,\dd K
=-\varphi'(1)(y-1)+\varphi(y)-\varphi(1).
\]
Substituting back into \eqref{eq:vspointwise} cancels everything but
$\varphi(y)$.  Now set $y=Y$ in \eqref{eq:vspointwise} and take
expectations.  The linear term carries $\E[Y]-1=0$ and dies; the two
kernels become put and call prices.
\end{proof}

In words: any smooth payoff is a constant, plus a forward position,
plus a strip of puts below the forward and calls above it, held in
amounts given by the payoff's second derivative.  The identity is
completely model-free --- it never asks where the prices came from.

Now specialize to the log contract, $\varphi(y)=-2\log y$.  Then
$\varphi(1)=0$ and $\varphi''(y)=2/y^{2}$, so
\eqref{eq:vsspanning} reads
\[
w_{\rm vs}
=2\int_0^{1}\frac{P(K)}{K^{2}}\,\dd K
+2\int_1^{\infty}\frac{c(K)}{K^{2}}\,\dd K .
\]
Every strike appears, weighted by $1/K^{2}$: the variance swap is a
strip of out-of-the-money options, densest below the forward.  For
computation it is tidier in log-moneyness.  Substitute $K=e^{k}$,
$\dd K=e^{k}\dd k$, and index the same prices by log-strike, writing
$c(k)$ for $c(e^{k})$ and likewise for the put; each $1/K^{2}\,\dd K$
becomes $e^{-k}\dd k$.  Now recall put--call parity in normalized
units ($P(k)=c(k)+(e^{k}-1)$, from \cref{eq:calldef}) and write
$(x)^-=\min(x,0)$ for the negative part.  On the call leg ($k>0$) the
quantity $e^{k}-1$ is positive, so $(e^{k}-1)^-=0$ and the bracket
below is just the call; on the put leg ($k<0$) it is negative, so
$(e^{k}-1)^-=e^{k}-1$ and the bracket is $c(k)+(e^{k}-1)=P(k)$,
exactly the put parity demands.  One bracket therefore covers both
legs, and the two integrals merge into one integral over the smile:
\begin{equation}
w_{\rm vs}
=2\int_{-\infty}^{\infty}
\Big[\Black\big(k,w(k)\big)+(e^{k}-1)^-\Big]\,e^{-k}\,\dd k ,
\label{eq:vsrepl}
\end{equation}
where $\Black$ is the Black call function of \cref{eq:black}, so
$\Black(k,w(k))=c(k)$ is the smile's call price.  Equation
\eqref{eq:vsrepl} consumes nothing but the total-variance curve
$w(k)$.  Any model that can state its smile can be integrated --- this
is the natural route for SVI and MCS, whose $w(k)$ is a formula.

\subsection{Against the law}

The second route never leaves the law.  The mean of any function of
$X$ can be taken in rank coordinates --- this is
\cref{eq:rankintegral}, the change of variables under which Chapter~2
built everything: if $U$ is uniform on $(0,1)$ then $Q(U)$ has the law
of $X$, so integrating the quantile function $Q$ over ranks
\emph{is} taking the mean.  Applying this to
$w_{\rm vs}=-2\,\E[X]$ gives Chapter~2's \cref{eq:varswap}:
\begin{equation*}
w_{\rm vs}
=-2\int_0^{1}Q(u)\,\dd u
=-2\int_{-\infty}^{\infty}x(z)\,\rho(z)\,\dd z ,
\end{equation*}
where the second form substitutes the log-odds coordinate $z$, with
$x(z)$ the transport and $\rho$ the logistic density
(\cref{sec:model}).  For an LQD slice --- one expiry's fitted law ---
this is one dot product against arrays the slice has already built on
the fixed log-odds grid every slice carries: no Black function, no
inversion, no new grid.  The integrand $-2\,x(z)\rho(z)$ decays like
$|z|e^{-|z|}$ --- Chapter~2's exponential tails make the transport
asymptotically linear in $z$, and the logistic weight supplies the
$e^{-|z|}$ --- so the slice's grid exhausts it.

\subsection{Along paths}

The third route uses neither prices nor the terminal law, but the
field.  Return to the left-hand side of \eqref{eq:vslogcontract},
the expected accumulated variance, in the representative model of
Chapter~4 where $v_s=v(s,Y_s)$.  The expectation of a time integral is
the time integral of expectations, and each instant's expectation is
an integral against that instant's density $f(s,y)$ --- the marginal
density of $Y_s$ that \cref{sec:lveq} evolved with the forward
equation:
\begin{equation}
w_{\rm vs}
=\E\!\left[\int_0^{\tau}v(s,Y_s)\,\dd s\right]
=\int_0^{\tau}\!\E\big[v(s,Y_s)\big]\,\dd s
=\int_0^{\tau}\!\!\int_0^{\infty}v(s,y)\,f(s,y)\,\dd y\,\dd s .
\label{eq:vsoccupation}
\end{equation}
Read the double integral as a pairing of two surfaces: the field
$v(s,y)$, against the sheet of densities $f(s,y)$ that says how much
time the paths spend at each level.  The density sheet is called the
\emph{occupation measure} of the diffusion, and \eqref{eq:vsoccupation}
says the fair strike is the field averaged over where the paths
actually go.

Computing \eqref{eq:vsoccupation} does not require assembling $f$.
Define the \emph{expected remaining variance}
\[
\mathcal{W}(s,y)
=\E\!\left[\int_s^{\tau}v(r,Y_r)\,\dd r\;\Big|\;Y_s=y\right],
\]
the fair strike of a swap started at time $s$ from level $y$.  It
satisfies a backward equation of exactly the form Chapter~4 marched,
with the field itself as a source term:
\begin{equation}
\partial_s\mathcal{W}
+\tfrac12\,v(s,y)\,y^{2}\,\partial_{yy}\mathcal{W}
+v(s,y)=0,
\qquad
\mathcal{W}(\tau,\cdot)=0,
\qquad
w_{\rm vs}=\mathcal{W}(0,1).
\label{eq:vssource}
\end{equation}
The derivation is two applications of ideas already in hand ---
It\^o's formula and the martingale property --- and is written out in
\cref{app:vsprotocol}.  Operationally: one backward sweep of the same
solver Chapter~4 marched forward, adding
$v\,\Delta s$ at each step, and the number is read off at the spot.
This is the natural route for the local-volatility model, which owns
neither a smile formula nor a terminal quantile --- but owns the
field.

\subsection{One number}

Collecting the three routes with the definition they all serve:
\begin{equation}
\boxed{\;
w_{\rm vs}
=-2\,\E[X]
=2\!\int_{\R}\!\Big[\Black\big(k,w(k)\big)+(e^{k}-1)^-\Big]e^{-k}\dd k
=-2\!\int_{\R}x(z)\,\rho(z)\,\dd z
=\int_0^{\tau}\!\!\int v\,f\,\dd y\,\dd s .
\;}
\label{eq:vsmaster}
\end{equation}
One number; an integral against the smile, an integral against the
law, an integral against the field.  Each equality is a theorem, so
any two routes evaluated on the same model must agree --- and
checking that they do is a free, end-to-end audit of an
implementation.  (Route two costs one \emph{quadrature} --- one
numerical integral --- and the audits below are limited only by
quadrature accuracy.)

\begin{figure}[htbp]
\centering
\paperfig{\textwidth}{fig_vs_three.pdf}
\caption{One number, three integrands.  (a)~The strike-side integrand
of \eqref{eq:vsrepl} on the running LQD slice; the orange ticks mark
the quoted strikes, and the shaded area is $w_{\rm vs}$.  (b)~The
law-side integrand $-2\,x(z)\rho(z)$ on the same slice's log-odds
grid; the signed area is the same number to \MacVsAgreeRankBp\ vol bp.
(c)~The field route on Chapter~4's calibrated SPY sheet: the local
volatility as a heatmap, the occupation measure as white density
contours funnelling out of the spot, and the marker at
$(y,\tau)=(1,0)$ where the backward solve \eqref{eq:vssource} reads
off the number.}
\label{fig:vsthree}
\end{figure}

\Cref{fig:vsthree} draws the three integrands for real fitted objects,
and it repays a slow look.  In panel~(a), the strike side on the
running LQD slice of \cref{fig:vspins}: the integrand of
\eqref{eq:vsrepl} peaks near the money --- the $e^{-k}$ weight decays,
but the option prices it multiplies decay much faster --- and the
orange ticks show how little of the strike axis the quotes occupy.
The area under this curve is the fair strike,
$\sigma_{\rm vs}=\MacVsThreeSvsPct\%$.  In panel~(b), the law side of
the \emph{same} slice: the integrand $-2\,x(z)\rho(z)$ against
log-odds.  It has two lobes, because $-2x(z)$ is positive on the
downside ranks ($x<0$) and negative on the upside: the green lobe is
downside outcomes pushing the fair strike up, the orange lobe upside
outcomes pulling it down.  The lobes are visibly unequal --- that
asymmetry is the skew --- and their \emph{signed} area agrees with
panel~(a)'s area to \MacVsAgreeRankBp\ vol bp: two utterly different
integrals, one number, the agreement limited only by quadrature.  In
panel~(c), the field route on a different model of the same market:
Chapter~4's SPY sheet, the field fitted to all eight expiries at
once.  The heatmap is the calibrated
local volatility; the white contours are the occupation measure ---
the densities $f(s,y)$ fanning out of the spot as time runs up the
axis --- and the pairing \eqref{eq:vsoccupation} weighs the heatmap by
the contours.  The marker at the origin of the fan is where
\eqref{eq:vssource} reads the answer.

The field route's audit deserves its own sentence, because it teaches
the same lesson as Chapter~4's two operators.  On the sheet, the
strike-side route \eqref{eq:vsrepl} and the field-side route
\eqref{eq:vssource} are both computed by finite differences, and both
carry a first-order bias in the time step.  At the march's own step
the two reads differ by \MacVsAgreeFieldCoarseBp\ vol bp; refine the
step eightfold and the gap collapses to \MacVsAgreeFieldBp\ vol bp,
halving with each halving of the step.  The mathematics says the
integrals are equal; the computation agrees --- in the limit, at the
first-order rate the schemes promise, and measuring that convergence
is the audit.

One structural observation organizes all of this, and it is worth
stating because it predicts which route each model finds cheap.  The
fair strike is \emph{linear} in the option prices and \emph{linear}
in the quantile $Q$ --- both appear under plain integral signs in
\eqref{eq:vsmaster} --- so a model whose own coordinates are one of
those objects gets the number, and by the same linearity its
parameter derivatives, for the price of one quadrature: LQD holds
$Q$, and the law route is a dot product.  It is \emph{not} linear in
the smile $w(k)$ --- the smile enters \eqref{eq:vsrepl} through the
nonlinear Black function --- so a model whose coordinates are the
smile must pay the Black evaluations: cheap when $w(k)$ is a formula,
as for SVI and MCS, and expensive precisely when the smile must be
\emph{solved for} point by point.  The field route's economy has a
third source.  Its pairing \eqref{eq:vsoccupation} is against an
occupation measure the field itself generates, so the number is not
linear in $v$.  What saves it is \eqref{eq:vssource}:
$\mathcal{W}$ solves a \emph{linear} equation in which $v$ enters
only as a coefficient and a source, so the number still costs one
backward sweep.  This section's rule, stated once: \emph{evaluate a
shared functional in each model's own chart}.  It stands behind
every route above.

We can now compute the number exactly, in any model, and trust the
result.  The next question is what part of it the quotes actually pin
down.
